\section{Phân tích sức khỏe và dân số}

\subsection{Vai trò trong bài toán chung}

Phân tích sức khỏe và dân số đánh giá mức độ phát triển ở khía cạnh chất lượng đời sống con người. Trong khung phân tích phát triển bền vững của nhóm, kinh tế phản ánh nguồn lực vật chất, giáo dục phản ánh năng lực con người, còn sức khỏe cho biết các nguồn lực đó có được chuyển hóa thành phúc lợi xã hội thực chất hay không. Vì vậy, các chỉ số như tuổi thọ trung bình, tử vong trẻ em dưới 5 tuổi, chi tiêu y tế và quy mô dân số được dùng để kiểm tra kết quả cuối cùng của quá trình phát triển.

Phần này cũng đóng vai trò kết nối với các phần trước: một quốc gia có GDP cao hoặc chi tiêu lớn chưa chắc đạt kết quả sức khỏe tốt nếu hệ thống y tế phân bổ nguồn lực không hiệu quả, tiếp cận dịch vụ không công bằng, hoặc chịu tác động mạnh từ các cú sốc như đại dịch, xung đột và khủng hoảng nhân đạo.

\subsection{Câu hỏi phân tích}

Phân tích sức khỏe và dân số tập trung vào ba câu hỏi chính:

\begin{enumerate}
    \item Tuổi thọ trung bình của các quốc gia đại diện thay đổi như thế nào trong giai đoạn 2000--2023? Đại dịch COVID-19 có tạo ra điểm gãy rõ rệt trong chuỗi thời gian không?
    \item Tỷ lệ tử vong trẻ em dưới 5 tuổi năm 2022 khác nhau như thế nào giữa các khu vực? Khu vực nào vẫn đối mặt với rủi ro sức khỏe trẻ em cao nhất?
    \item Chi tiêu y tế theo \% GDP có quan hệ như thế nào với tuổi thọ trung bình? Phân bố tuổi thọ theo khu vực cho thấy mức độ bất bình đẳng sức khỏe ra sao?
\end{enumerate}

\subsection{Chỉ số sử dụng (Indicators)}

Bốn chỉ số WDI được chọn cho mục tiêu sức khỏe và dân số, bao gồm hai chỉ số kết quả sức khỏe, một chỉ số đầu vào hệ thống y tế và một chỉ số quy mô dân số.

\begin{longtable}{L{3.8cm}L{3.5cm}L{1.3cm}L{5cm}}
\caption{Các chỉ số WDI sử dụng trong phân tích sức khỏe và dân số, kèm theo vai trò phân tích.}
\label{tab:health-indicators}\\
\toprule
Mã WDI & Tên cột output & Đơn vị & Ý nghĩa và lý do chọn \\
\midrule
\endfirsthead
\toprule
Mã WDI & Tên cột output & Đơn vị & Ý nghĩa và lý do chọn \\
\midrule
\endhead
\path{SP.DYN.LE00.IN} & \path{Life_expectancy} & năm & Tuổi thọ trung bình tính từ lúc sinh; là thước đo outcome tổng hợp phản ánh điều kiện sống, chất lượng y tế và mức độ phát triển xã hội. \\
\path{SH.DYN.MORT} & \path{Under5_mortality} & trên 1.000 trẻ sinh sống & Tỷ lệ tử vong trẻ em dưới 5 tuổi; phản ánh chất lượng chăm sóc sức khỏe bà mẹ--trẻ em, dinh dưỡng, tiêm chủng và vệ sinh cơ bản. \\
\path{SH.XPD.CHEX.GD.ZS} & \path{Health_expenditure_pct_GDP} & \% GDP & Chi tiêu y tế hiện thời; đo mức nguồn lực đầu vào mà quốc gia dành cho hệ thống y tế. \\
\path{SP.POP.TOTL} & \path{Population} & người & Tổng dân số; dùng để đặt các kết quả sức khỏe trong bối cảnh quy mô xã hội, đồng thời làm biến kích thước trong biểu đồ bubble. \\
\bottomrule
\end{longtable}

\subsection{Tiền xử lý dữ liệu}

Dữ liệu sức khỏe được trích từ file \texttt{WDIEXCEL.xlsx} theo pipeline tiền xử lý chung của nhóm. Quy trình cụ thể:

\begin{itemize}
    \item Lọc 4 indicators sức khỏe--dân số trong giai đoạn 2000--2023.
    \item Chỉ giữ các quốc gia hoặc vùng lãnh thổ có \texttt{Region} hợp lệ trong sheet \texttt{Country}; các dòng tổng hợp như World, nhóm thu nhập hoặc khu vực tổng hợp bị loại bỏ.
    \item Với từng cặp \texttt{Country\_Code}--\texttt{Indicator\_Code}, loại chuỗi không có quan sát hoặc có khoảng thiếu nội bộ dài hơn 2 năm liên tiếp.
    \item Nội suy tuyến tính các khoảng thiếu nội bộ dài tối đa 2 năm; không ngoại suy ở đầu hoặc cuối chuỗi thời gian.
    \item Pivot dữ liệu từ file sạch chung sang wide format và xuất file \path{data_output/wdi_health.csv}.
\end{itemize}

Sau tiền xử lý, file \path{wdi_health.csv} có \textbf{5.208 dòng} và \textbf{217 quốc gia/vùng lãnh thổ}. Hai chỉ số \path{Life_expectancy} và \path{Population} có độ phủ đầy đủ cho toàn bộ 217 quốc gia trong giai đoạn 2000--2023. Hai chỉ số \path{Under5_mortality} và \path{Health_expenditure_pct_GDP} có mức thiếu cao hơn, nhưng vẫn đủ rộng để phân tích so sánh khu vực.

\begin{longtable}{L{3.8cm}L{2.5cm}L{2.5cm}L{2.5cm}L{2.5cm}}
\caption{Mức phủ dữ liệu của 4 indicator sức khỏe sau tiền xử lý.}
\label{tab:health-coverage}\\
\toprule
Indicator & Quốc gia giữ & Quốc gia loại (rỗng) & Quốc gia loại (gap) & Tổng bản ghi \\
\midrule
\endfirsthead
\toprule
Indicator & Quốc gia giữ & Quốc gia loại (rỗng) & Quốc gia loại (gap) & Tổng bản ghi \\
\midrule
\endhead
\path{SP.DYN.LE00.IN} & 217 & 0 & 0 & 5.208 \\
\path{SH.DYN.MORT} & 196 & 21 & 0 & 4.704 \\
\path{SH.XPD.CHEX.GD.ZS} & 193 & 24 & 0 & 4.540 \\
\path{SP.POP.TOTL} & 217 & 0 & 0 & 5.208 \\
\bottomrule
\end{longtable}

\subsection{Trực quan hóa và nhận xét}

% ======================================================================
\subsubsection{Biểu đồ 1 --- Xu hướng tuổi thọ trung bình giai đoạn 2000--2023 (Line chart)}

\textbf{Loại biểu đồ:} Biểu đồ đường nhiều chuỗi (multi-line chart).\\
\textbf{Lý do chọn:} Tuổi thọ là chỉ số liên tục theo thời gian, do đó biểu đồ đường phù hợp để quan sát xu hướng dài hạn, nhận diện điểm gãy trong chuỗi thời gian và so sánh các quốc gia đại diện trên cùng hệ trục.

\reportfigure{images/health/health_life_expectancy_trend.png}{Xu hướng tuổi thọ trung bình của các quốc gia đại diện trong giai đoạn 2000--2023. Trục tung biểu diễn tuổi thọ trung bình tính theo năm; mỗi đường thể hiện một quốc gia trong nhóm so sánh.}{fig:health-le-trend}

\textbf{Nhận xét:}
\begin{itemize}
    \item \textbf{Xu hướng cải thiện dài hạn:} Hầu hết các đường đều tăng từ năm 2000 đến trước năm 2020. Nhóm đường cao nhất nằm quanh 80--84 năm, nhóm trung gian tăng dần từ khoảng 63--73 năm lên vùng 72--78 năm, còn đường thấp nhất tăng từ khoảng 47 năm lên gần 55 năm.
    \item \textbf{Điểm gãy rõ trên hình:} Một số đường giảm xuống ở năm 2021 trước khi phục hồi sau đó. Nhịp giảm thấy rõ nhất ở các đường đang nằm trong vùng 70--80 năm, cho thấy cú sốc sức khỏe giai đoạn đại dịch làm gián đoạn xu hướng tăng đều trước đó.
    \item \textbf{Khoảng cách sức khỏe vẫn lớn:} Đến năm 2023, đường cao nhất vẫn ở trên 80 năm, trong khi đường thấp nhất chỉ quanh 54--55 năm. Chênh lệch xấp xỉ 30 năm trên biểu đồ cho thấy bất bình đẳng tuổi thọ giữa các nhóm quốc gia vẫn rất đáng kể.
\end{itemize}

% ======================================================================
\subsubsection{Biểu đồ 2 --- Tử vong trẻ em dưới 5 tuổi theo khu vực năm 2022 (Bar chart)}

\textbf{Loại biểu đồ:} Biểu đồ cột ngang (horizontal bar chart).\\
\textbf{Lý do chọn:} Biểu đồ cột ngang phù hợp để so sánh một chỉ số định lượng giữa các nhóm phân loại, đặc biệt khi tên khu vực dài và cần sắp xếp thứ hạng rõ ràng.

\reportfigure{images/health/health_under5_mortality_bar.png}{Tỷ lệ tử vong trẻ em dưới 5 tuổi trung bình theo khu vực năm 2022. Đơn vị là số trẻ tử vong trên 1.000 trẻ sinh sống.}{fig:health-mortality-bar}

\textbf{Nhận xét:}
\begin{itemize}
    \item \textbf{Sub-Saharan Africa là khu vực rủi ro cao nhất:} Thanh của khu vực này dài vượt trội, xấp xỉ 60--70 ca trên 1.000 trẻ sinh sống theo trục hoành. Mức này cao hơn nhiều lần so với tất cả khu vực còn lại, cho thấy gánh nặng tử vong trẻ em tập trung rất mạnh ở khu vực này.
    \item \textbf{Nhóm trung bình:} South Asia, Middle East \& North Africa và East Asia \& Pacific có chiều dài thanh khá gần nhau, đều quanh vùng 20 ca trên 1.000. Latin America \& Caribbean thấp hơn một chút, khoảng 15--17 ca.
    \item \textbf{Nhóm thấp nhất là Europe \& Central Asia và North America:} Hai thanh cuối ngắn nhất, chỉ khoảng 6--7 ca trên 1.000. Khoảng cách giữa nhóm này và Sub-Saharan Africa rất lớn, phản ánh sự khác biệt rõ về điều kiện chăm sóc sức khỏe trẻ em.
\end{itemize}

% ======================================================================
\subsubsection{Biểu đồ 3 --- Chi tiêu y tế và tuổi thọ năm 2022 (Bubble scatter plot)}

\textbf{Loại biểu đồ:} Biểu đồ phân tán kết hợp kích thước bong bóng (bubble scatter plot).\\
\textbf{Lý do chọn:} Scatter plot phù hợp để kiểm tra mối liên hệ giữa hai biến định lượng: chi tiêu y tế theo \% GDP và tuổi thọ trung bình. Kích thước bong bóng bổ sung chiều thông tin thứ ba là quy mô dân số, giúp các quốc gia đông dân được đặt trong đúng bối cảnh tác động xã hội.

\reportfigure{images/health/health_expenditure_expectancy_scatter.png}{Mối liên hệ giữa chi tiêu y tế (\% GDP) và tuổi thọ trung bình năm 2022. Mỗi điểm là một quốc gia; kích thước bong bóng phản ánh dân số và màu sắc thể hiện khu vực.}{fig:health-bubble-scatter}

\textbf{Nhận xét:}
\begin{itemize}
    \item \textbf{Mối quan hệ không tuyến tính rõ ràng:} Các điểm tập trung dày nhất trong vùng chi tiêu khoảng 3--12\% GDP và tuổi thọ khoảng 65--85 năm. Không có một đường xu hướng tăng thật rõ, nên không thể kết luận từ hình rằng chi tiêu y tế cao hơn luôn đi kèm tuổi thọ cao hơn.
    \item \textbf{Nhiều quốc gia đạt tuổi thọ cao ở mức chi tiêu trung bình:} Vùng 5--10\% GDP có nhiều điểm nằm trên 75 năm, cho thấy hiệu quả sức khỏe không chỉ phụ thuộc vào tỷ trọng chi tiêu mà còn liên quan đến cách tổ chức hệ thống y tế và điều kiện xã hội.
    \item \textbf{Có các điểm ngoại lệ đáng chú ý:} Một số điểm ở bên phải biểu đồ có chi tiêu trên 16\% GDP nhưng tuổi thọ chỉ quanh 65--78 năm. Ngoài ra, có một điểm rất thấp gần 19 năm tuổi thọ ở khoảng 10\% GDP, làm nổi bật khả năng tồn tại các trường hợp bất thường hoặc chịu tác động khủng hoảng nghiêm trọng.
    \item \textbf{Giá trị null cần được lưu ý:} Hình thể hiện 25 nulls ở góc dưới bên phải, cho thấy không phải mọi bản ghi đều có đủ dữ liệu để vẽ. Vì vậy, nhận xét từ scatter plot nên được hiểu là mô tả trên phần dữ liệu hiển thị.
\end{itemize}

% ======================================================================
\subsubsection{Biểu đồ 4 --- Phân bố tuổi thọ theo khu vực năm 2022 (Box plot)}

\textbf{Loại biểu đồ:} Biểu đồ hộp (box plot).\\
\textbf{Lý do chọn:} Box plot phù hợp để so sánh phân bố của một biến định lượng giữa các nhóm. Biểu đồ không chỉ cho thấy trung vị, mà còn thể hiện độ phân tán, khoảng tứ phân vị và các giá trị ngoại lai trong từng khu vực.

\reportfigure{images/health/health_life_expectancy_boxplot.png}{Phân bố tuổi thọ trung bình theo khu vực năm 2022. Mỗi hộp thể hiện phân bố tuổi thọ của các quốc gia trong cùng một khu vực.}{fig:health-le-boxplot}

\textbf{Nhận xét:}
\begin{itemize}
    \item \textbf{North America có phân bố cao và hẹp nhất:} Hộp của North America nằm chủ yếu quanh 80--82 năm và khá nhỏ, cho thấy tuổi thọ giữa các quốc gia trong nhóm này tương đối đồng đều.
    \item \textbf{Europe \& Central Asia và East Asia \& Pacific có độ phân tán rộng hơn:} Hai nhóm này có nhiều điểm trải từ khoảng 60--85 năm, trong đó phần hộp nằm ở vùng 70--80 năm. Điều này cho thấy trong cùng khu vực vẫn có chênh lệch đáng kể giữa các quốc gia.
    \item \textbf{Sub-Saharan Africa thấp nhất và phân tán mạnh:} Hộp của khu vực này nằm thấp hơn các nhóm khác, chủ yếu quanh 60--67 năm. Ngoài ra còn có một điểm ngoại lệ rất thấp gần 19 năm, làm kéo sự chú ý đến các rủi ro sức khỏe cực đoan.
    \item \textbf{Các khu vực còn lại nằm ở nhóm trung gian:} Latin America \& Caribbean, Middle East \& North Africa và South Asia nằm giữa nhóm cao nhất và Sub-Saharan Africa. Nhìn từ biểu đồ, tuổi thọ toàn cầu không chỉ khác nhau giữa khu vực mà còn khác đáng kể ngay trong từng khu vực.
\end{itemize}

\subsection{Thảo luận}

Tổng hợp bốn biểu đồ cho thấy bức tranh sức khỏe toàn cầu có cả tiến bộ rõ rệt lẫn bất bình đẳng sâu sắc:

\begin{itemize}
    \item \textbf{Tuổi thọ tăng nhưng không đồng đều:} Giai đoạn 2000--2019 cho thấy xu hướng tăng tuổi thọ ở phần lớn quốc gia, nhưng khoảng cách giữa nhóm phát triển và nhóm thu nhập thấp vẫn rất lớn. Tăng trưởng kinh tế và mở rộng dịch vụ y tế có tác động tích cực, nhưng mức xuất phát khác nhau khiến chênh lệch tuyệt đối vẫn kéo dài.
    \item \textbf{Sức khỏe trẻ em là điểm phân hóa mạnh nhất:} Tử vong trẻ em dưới 5 tuổi ở Sub-Saharan Africa cao hơn nhiều lần so với North America và Europe \& Central Asia. Đây là chỉ báo quan trọng vì tử vong trẻ em thường gắn với các điều kiện nền tảng như dinh dưỡng, tiêm chủng, nước sạch, vệ sinh và chăm sóc y tế ban đầu.
    \item \textbf{Chi tiêu y tế không phải yếu tố quyết định duy nhất:} Scatter plot cho thấy quan hệ giữa chi tiêu y tế theo \% GDP và tuổi thọ không tuyến tính rõ. Một quốc gia có thể chi tiêu nhiều nhưng kết quả tuổi thọ không tương xứng nếu hệ thống y tế kém hiệu quả hoặc bất bình đẳng tiếp cận dịch vụ cao. Ngược lại, một số quốc gia đạt tuổi thọ cao với mức chi tiêu vừa phải nhờ y tế dự phòng, quản trị tốt và điều kiện xã hội thuận lợi.
    \item \textbf{Khả năng chống chịu trước cú sốc là vấn đề trọng tâm:} Điểm gãy 2020--2021 trong biểu đồ tuổi thọ cho thấy đại dịch COVID-19 làm suy giảm outcome sức khỏe ở nhiều quốc gia. Điều này nhấn mạnh rằng phát triển sức khỏe bền vững không chỉ là cải thiện trung bình dài hạn, mà còn là năng lực duy trì hệ thống y tế trước khủng hoảng.
\end{itemize}

Khi đặt cạnh phần kinh tế và giáo dục, kết quả sức khỏe cho thấy phát triển bền vững cần được hiểu như một chuỗi liên kết: nguồn lực kinh tế tạo điều kiện đầu tư, giáo dục nâng cao năng lực xã hội, còn hệ thống y tế quyết định khả năng chuyển hóa các nguồn lực đó thành tuổi thọ cao hơn, tử vong trẻ em thấp hơn và chất lượng sống tốt hơn.

\subsection{Kết luận mục tiêu sức khỏe và dân số}

Phân tích sức khỏe và dân số cho thấy thế giới đã đạt nhiều tiến bộ trong hơn hai thập kỷ qua: tuổi thọ nhìn chung tăng, tử vong trẻ em giảm, và nhiều quốc gia đang phát triển cải thiện nhanh chóng. Tuy nhiên, bất bình đẳng sức khỏe vẫn là vấn đề lớn. Sub-Saharan Africa tiếp tục là khu vực chịu gánh nặng cao nhất về tử vong trẻ em và tuổi thọ thấp, trong khi các khu vực phát triển duy trì tuổi thọ cao và phân bố ổn định hơn.

Kết quả cũng cho thấy không thể đánh giá hệ thống y tế chỉ bằng mức chi tiêu. Chi tiêu y tế là điều kiện cần, nhưng hiệu quả phân bổ, y tế dự phòng, công bằng tiếp cận, giáo dục, thu nhập và ổn định xã hội mới quyết định khả năng chuyển hóa nguồn lực thành kết quả sức khỏe thực chất. Đây là điểm then chốt khi nhóm tổng hợp bốn mục tiêu để đánh giá phát triển bền vững trong toàn bài.
